% ============================================================================
% Chapter 7 --- Policy decision and enforcement
% ----------------------------------------------------------------------------
% Purpose : PDP/PEP layer; policy-as-code; decision logging.
%           Belongs to the Control and governance plane (D5.1).
% Page target: ~5 pages.
% Dependencies: Ch. 5 (standards: OPA REST, Cedar API), Ch. 6 (identity
%               attributes consumed as policy input), Ch. 3 (sovereignty
%               framework).
% Mirrors the layer chapter template established in Ch. 6 (D6.5).
% ----------------------------------------------------------------------------
% Decisions registered in _coherence-EN.md:
%   D7.1 minimum compliance set for a policy decision point (the contract)
%   D7.2 reference instantiation = OPA + bundle server + decision log
%   D7.3 alternatives are Cedar (open formal semantics) and native IdP authz
%   D7.4 sovereignty scores for the policy layer (NEW, not in Ch. 3)
%   D7.5 the policy plane is the most ambitious novelty for many adopters
% ============================================================================

\chapter{Policy decision and enforcement}
\label{ch:policy}

\begin{caveatbox}[Dependencies]
\noindent This chapter applies the layer template of
\trchap{ch:identity}{} to the policy layer. It consumes identity
attributes from \trchap{ch:identity}{} and emits decisions consumed
by Network (\trchap{ch:network}{}), Compute
(\trchap{ch:compute}{}), Storage (\trchap{ch:storage}{}) and
Endpoint (\trchap{ch:endpoint}{}). The decision log produced by this
layer is consumed by Observability (\trchap{ch:observability}{}).
The layer belongs to the \emph{Control and governance plane}.
\end{caveatbox}

The policy layer answers the second question of any access decision,
once the identity layer has answered the first: \emph{given who is
acting, with which attributes and from which posture, what is
permitted?}. It is the layer that translates institutional rules ---
data classification, role-activation gating, time-of-day or
location-of-access conditions, posture requirements --- into
machine-readable decisions consumed by the execution plane. It is
also, for many adopters at the time of writing --- particularly
those without a dedicated platform-engineering function --- the most
ambitious novelty in this architecture: relatively few institutions
today operate a dedicated policy decision point separate from their
applications and identity provider. Adopting one is a real
organisational change, addressed in \trchap{ch:wave4-policy}{Wave~4
--- Policy plane}.

%-----------------------------------------------------------------------------
\section{Functional role and interface contract}
\label{sec:po-contract}

The policy layer is responsible for: receiving a decision request
that combines a subject (with attributes and posture claims), a
resource (with classification and ownership), and an environment
(with time, location, and contextual signals); evaluating that
request against the institution's policy corpus; and returning a
structured decision that downstream enforcement points consume. It
is also responsible for emitting an audit-grade decision log that
makes every consequential decision reconstructable.

\begin{contractbox}[Policy layer]
\noindent A compliant policy decision point, in the sense of this
architecture, must:

\begin{itemize}[leftmargin=1.5em,nosep]
  \item expose a query API, either OPA REST or Cedar evaluation API,
    receiving a JSON-structured request and returning a JSON-structured
    decision;
  \item evaluate against a policy corpus expressed in a high-level
    language with formal evaluation semantics (Rego, Cedar);
  \item support a documented \emph{bundle distribution} pattern by
    which policies are versioned, signed, and pushed to the decision
    point without redeployment;
  \item return decisions of the form
    \texttt{\{allow, deny, indeterminate\}} accompanied by structured
    \emph{obligations} (e.g.\ session timeout, audit class, additional
    factor required) that the enforcement point applies;
  \item produce a decision log that includes the inputs, the policy
    version, the decision and the obligations, in a form ingestible
    by the observability layer;
  \item support a \emph{shadow mode} in which the decision is logged
    but not enforced, used during introduction or after a policy
    change.
\end{itemize}
\end{contractbox}

The contract is silent on the policy authoring workflow (who writes
policies, how they are reviewed, how they are deployed) and on the
mapping of institutional rules to the policy language. These are
governance choices addressed in \trchap{ch:wave4-policy}{}; the
contract requires only that the chosen workflow respect the
versioning and the audit-trail properties.

%-----------------------------------------------------------------------------
\section{Reference instantiation: OPA + bundle server + decision log to SIEM}
\label{sec:po-reference}

The reference instantiation comprises an Open Policy Agent
(OPA)~\cite{opa-project}
deployment as the decision point, a bundle server as the policy
distribution mechanism, and a decision-log forwarder that delivers
every decision to the observability pipeline.

\begin{instantiationbox}[Policy layer]
\noindent\textbf{Topology.} OPA runs as a sidecar or as a centralised
service depending on latency requirements; policies are authored in
Rego and stored in a Git repository under standard code review;
\emph{bundles} (signed archives of compiled policies and data) are
built by continuous integration and served by a bundle server, from
which OPA instances pull at a configurable interval; decision logs
are emitted in OpenTelemetry format and routed to the observability
layer (\trchap{ch:observability}{}). The bundle distribution pattern
is co-normative with OPA REST itself: it is the mechanism that turns
a runtime decision point into a governable artefact.
\end{instantiationbox}

\paragraph{Authoring workflow.} Policies are authored in Rego in a
Git repository under code review. Each merge produces a versioned
bundle, signed by the build system, served by the bundle server.
\gls{opa} instances pull the bundle and reload without restart. The
review workflow is the institutional governance mechanism for the
policy corpus; the bundle distribution is the technical mechanism
that propagates the result.

\paragraph{Operational considerations.} The reference instantiation
is expected to require expertise in Rego (or in the Cedar
alternative, see below), which is a domain-specific language whose
learning curve is typically reported in weeks rather than days. The
principal recurring effort is policy testing: each policy is
unit-tested against representative request patterns, and the test
suite is maintained alongside the policy corpus. The operational
team would typically be small --- indicatively a fraction of the
institution's platform-engineering staffing profile, shared across
other duties --- but the
institutional skill investment remains significant.

\paragraph{Decision log schema.} The schema includes, at minimum:
the request inputs (subject identifier, resource identifier,
environment signals), the policy version, the decision verdict, the
obligations, the evaluation duration, and a correlation identifier
that ties the decision to the activity context. The canonical schema
is consolidated in \trchap{ch:observability}{} \S\ref{sec:ob-schema};
the policy layer is its principal producer.

\begin{realworldbox}[Policy plane in academic and research settings]
\noindent The policy plane in academia is more thinly documented
in policy-as-code terms than the identity plane. Authorisation
patterns relevant to research collaborations are developed in
GÉANT's eduTEAMS~\cite{eduteams} service, which implements the
AARC Blueprint Architecture~\cite{aarc-blueprint} --- a
sector-specific design pattern for
authentication and authorisation in international research
collaborations. The AARC Blueprint and policy-as-code engines
such as OPA address overlapping concerns at different levels:
the former is a sector-specific architecture, the latter a
general decision engine; their combination remains a relatively
open area for academic deployments. Industrial OPA deployments
publicly described by large-scale technology and financial-services
organisations demonstrate the engine at scale; their
transposition to academic settings is conditional on the
institution's authorisation profile. The phase-2 international
university network anticipated by this Reference would be one of
the contexts in which a shared library of academic-context
policies might emerge.
\end{realworldbox}

%-----------------------------------------------------------------------------
\section{Alternative~1 --- Cedar with the Verified Permissions pattern}
\label{sec:po-alt1}

\begin{alternativebox}[Cedar evaluation engine]
\noindent Cedar~\cite{cedar-project} is an open-source policy language,
developed by AWS and released as open source, designed for static
verifiability. It
admits a formal semantics that supports automated checks (for
example, that no policy in the corpus grants a forbidden action under
any input). Cedar is integrated either standalone or as part of the
AWS Verified Permissions service; the architecture relies on Cedar
itself, not on the AWS hosted variant, to preserve the technical
sovereignty dimension.
\end{alternativebox}

\paragraph{Where it adds value.} Cedar's formal semantics make the
policy corpus more readable than Rego in many cases, and they admit
verification techniques (consistency, no-bypass) that Rego does not.
For institutions whose risk posture demands a higher assurance on the
policy corpus itself, Cedar is the appropriate choice.

\paragraph{Where it constrains the institution.} The Cedar
ecosystem is younger than the OPA ecosystem; integrations with
specific enforcement points (network attachment, Kubernetes
admission, application middlewares) are fewer and less mature. The
institutional skill investment is comparable to Rego, but the labour
market is shallower at the time of writing.

\paragraph{When it is the right choice.} An institution with a
strong formal-methods culture, with a policy corpus large enough to
benefit from verifiability, and with the resources to bridge gaps in
the Cedar ecosystem with internal engineering. Cedar is also a
reasonable choice for an institution that wants to avoid OPA's
specific operational patterns (sidecar density, bundle refresh) and
prefers a request-response engine.

%-----------------------------------------------------------------------------
\section{Alternative~2 --- Native identity-provider authorisation rules}
\label{sec:po-alt2}

\begin{alternativebox}[Native IdP authorisation]
\noindent The institution does not introduce a separate policy
decision point. Authorisation rules are expressed in the identity
provider's own configuration: conditional-access expressions for
the institution's incumbent proprietary identity suite, authorisation
policies and scopes for Keycloak.
The identity layer is also the policy layer.
\end{alternativebox}

\paragraph{Where it works.} For institutions whose authorisation
needs are limited to coarse-grained role gating, MFA enforcement and
basic time-of-day or location conditions, the identity provider's
native rules are sufficient. The reduction in moving parts is real,
and the operational burden is lower.

\paragraph{Where it fails.} TAC-aware decisions require composing
multiple signals --- subject attributes, posture, resource
classification, time, location, environment --- in expressions whose
complexity quickly exceeds what a proprietary conditional-access
language accommodates without becoming unmaintainable. More fundamentally,
encoding the policy corpus in the identity provider couples it to
that provider: an exit from the identity provider becomes an exit
from the policy plane as well. Native authorisation is a viable
\emph{starting} configuration but a poor \emph{end-state} for an
institution whose sovereignty posture aims at separation of
concerns.

\paragraph{When it is acceptable.} As an early stage of the
migration, before the policy plane is mature enough to take over;
or as a permanent solution for institutions whose authorisation
needs do not justify the cost of a separate policy layer.

%-----------------------------------------------------------------------------
\section{Sovereignty grid applied}
\label{sec:po-sovereignty}

The three instantiations are profiled below using the grid of
\trchap{ch:sovereignty-grid}{}. The scores are specific to the
policy layer and are introduced in this chapter; they were not
worked through in \trchap{ch:sovereignty-grid}{}.

\begin{table}[H]
\centering\small
\caption{Per-dimension sovereignty profile of the three policy
instantiations.}
\label{tab:po-sovereignty}
\begin{tabularx}{\textwidth}{%
  >{\hsize=0.18\hsize\linewidth=\hsize\bfseries\raggedright\arraybackslash}X%
  >{\hsize=0.40\hsize\linewidth=\hsize\centering\arraybackslash}X%
  >{\hsize=0.40\hsize\linewidth=\hsize\centering\arraybackslash}X%
  >{\hsize=0.40\hsize\linewidth=\hsize\centering\arraybackslash}X%
  >{\hsize=1.62\hsize\linewidth=\hsize\raggedright\arraybackslash}X%
}
\toprule
Dimension & OPA + bundle & Cedar & Native IdP authz \\
\midrule
Data          & 3 Robust      & 3 Robust      & 1--2 (per IdP) \\
Technical     & 3 Robust      & 3 Robust      & 1 Partial \\
Financial     & 3 Robust      & 3 Robust      & inherits IdP \\
Operational   & 2 Established & 1--2 Partial  & 3 Robust \\
Institutional & 3 Robust      & 3 Robust      & 1 Partial \\
\bottomrule
\end{tabularx}
\end{table}

\noindent OPA and Cedar score similarly on the four sovereignty
dimensions other than operational, where the OPA labour market is
deeper. The native IdP path inherits the IdP's profile on most
dimensions and is consequently bounded by it: this is a structural
property of co-locating identity and policy, not a fixable gap. A
shift from native IdP authorisation to a separate policy plane
(OPA or Cedar) is consequently the highest-leverage sovereignty
move available in the control plane.

%-----------------------------------------------------------------------------
\section{Regulatory anchoring}
\label{sec:po-regulatory}

\noindent The policy layer is the per-request decision point that
operationalises the applicable data-protection regime's
data-protection-by-design obligation and contributes to its
access-control-effectiveness obligation. Under the applicable
cybersecurity-baseline regime --- its risk-analysis and
information-system security-policy requirement and its
access-control-policy requirement --- it is the locus where
the institutional risk posture is enforced uniformly. Under the
applicable AI regulation, its prohibited-practices and
transparency provisions
admit a policy-plane enforcement: rules that gate AI-related TACs
on context labelling and on prohibited-use classification. The
layer operationalises 27001~A.5.15 (access control), A.5.18
(access rights) and A.8.3 (information access restriction) and
contributes to NIST CSF \textsf{PROTECT} (PR.AA). Detailed mapping
in \trchap{ch:regulatory-mapping}{} \S\ref{sec:rm-gdpr},
\S\ref{sec:rm-nis2}, \S\ref{sec:rm-ai-act}, \S\ref{sec:rm-iso}.

%-----------------------------------------------------------------------------
\section{Common pitfalls}
\label{sec:po-pitfalls}

\begin{pitfallbox}[Policy explosion]
\noindent A policy corpus left ungoverned grows by accretion: every
new application adds rules, no rules are retired, and the corpus
loses the structure that makes it auditable. The mitigation is a
governance process for the corpus that mirrors the governance process
for application code: review, versioning, retirement, and a
pre-merge consistency check that the new rule does not conflict with
an existing one. Policy languages with formal semantics (Cedar in
particular) support automated consistency checks; OPA admits the same
through unit tests against representative inputs.
\end{pitfallbox}

\begin{pitfallbox}[Missing shadow mode]
\noindent A policy change deployed directly in enforcing mode
without a shadow-mode period silently denies legitimate access until
the affected users open tickets. The mitigation is to make shadow
mode the default for any new policy and any modification to an
existing policy: the decision is logged, compared against the
incumbent rule, and switched to enforcing only when the discrepancy
profile is understood. The shadow / enforcing toggle is part of the
contract of \S\ref{sec:po-contract}.
\end{pitfallbox}

\begin{pitfallbox}[Decision/log drift]
\noindent The decision log is intended to allow reconstruction of any
consequential decision, but a poorly-implemented log records a
summary that does not allow reconstruction (e.g.\ the obligations are
elided, the policy version is not recorded). The mitigation is to
exercise the reconstruction in the conformance suite of
\trchap{ch:conformance}{}: take a decision from the log and replay
it against the recorded policy version; the result must match.
\end{pitfallbox}

\begin{pitfallbox}[Hidden coupling to identity-provider attributes]
\noindent Policies that reference identity-provider-specific
attributes (e.g.\ vendor-specific extension claims from a
proprietary identity suite) lose portability
when the identity provider is replaced. The mitigation is to map the
identity-provider attributes to a stable institutional schema before
they reach the policy plane, and to author policies against the
institutional schema only. The mapping is itself part of the
contract between Identity and Policy.
\end{pitfallbox}

\begin{pitfallbox}[Over-centralisation of policy authoring]
\noindent A central policy team that becomes the bottleneck for every
application's authorisation rules slows the institution to a halt
and accumulates a corpus that the team alone understands. The
mitigation is to delegate authoring rights with code review:
application owners propose policies for their resources; the central
team reviews for consistency and conflicts. The model mirrors the
federated-governance delegation pattern common in mature
higher-education and research IT organisations.
\end{pitfallbox}
